\documentclass[11pt]{article}

\usepackage[margin=1in]{geometry}
\usepackage[T1]{fontenc}
\usepackage{amsmath,amssymb,amsthm}
\usepackage{mathtools}
\usepackage{stmaryrd}
\usepackage[protrusion=true,expansion=false]{microtype}
\usepackage[colorlinks=true,linkcolor=black,citecolor=black,urlcolor=blue]{hyperref}
\usepackage{enumitem}

% ---- theorem environments ----
\theoremstyle{definition}
\newtheorem{definition}{Definition}[section]
\newtheorem{example}[definition]{Example}
\theoremstyle{plain}
\newtheorem{theorem}[definition]{Theorem}
\newtheorem{proposition}[definition]{Proposition}
\newtheorem{lemma}[definition]{Lemma}
\newtheorem{corollary}[definition]{Corollary}
\newtheorem{conjecture}[definition]{Conjecture}
\theoremstyle{remark}
\newtheorem{remark}[definition]{Remark}

% ---- notation ----
\newcommand{\Sig}{\mathsf{Sig}}
\newcommand{\Stk}{\mathsf{Stk}}
\newcommand{\Hist}{\mathsf{Hist}}
\newcommand{\Ev}{\mathsf{Ev}}
\newcommand{\Real}{\mathbb{R}}
\newcommand{\Rpos}{\mathbb{R}_{>0}}
\newcommand{\cost}{\mathsf{C}}
\newcommand{\hist}{\mathsf{H}}
\newcommand{\ciGSLT}{\mathsf{ciGSLT}}
\newcommand{\Grph}{\mathcal{G}}
\newcommand{\cover}{\widetilde{\mathcal{G}}}
\newcommand{\im}{\operatorname{im}}
\newcommand{\id}{\mathrm{id}}
\newcommand{\fund}{\pi_1}
\newcommand{\emptyhist}{\varepsilon}
\DeclareMathOperator{\Betti}{b_1}
\DeclareMathOperator{\hol}{hol}

\title{\bfseries History as an Endofunctor\\[0.35em]
\large Reversibility, Covering Spaces, and the Memory Price of Energy Conservation\\ in Cost-Accounted GSLTs}
\author{L.~G.~Meredith\\ \small F1R3FLY.io}
\date{July 2026 --- working note}

\begin{document}
\maketitle

\begin{abstract}
Keeping a history as part of the state, and appending to it at every rewrite, recovers
reversibility: a rewrite theory that discards its past is irreversible precisely because a
configuration may have many predecessors, and recording which step was taken restores a unique
one. We package this as an endofunctor $\hist$ on the category $\ciGSLT$ of continued interactive
graph-structured lambda theories---the writer monad over the free monoid of rewrite events---and
observe that it is the time-reversal dual of the cost endofunctor $\cost$: where $\cost$ debits a
forward fuel modulus that \emph{drains} with reduction, $\hist$ accrues a backward trace that
\emph{accumulates}, the two being the ledgers of what remains and what was spent along one time
axis. The free history makes the graph of computation paths out of a term into a \emph{tree}: it
is the universal cover of the reduction graph, and reversibility is exactly its unique-parent
property. Erasures on the history form a lattice of monad quotients between this reversible cover
and the irreversible base; the injective ones compress the log without folding the tree, while the
identifying ones fold it, closing loops so that paths revisit states---an intermediate cover, i.e.\
a quotient of the reduction graph's fundamental group by a chosen set of loops, with the
order-forgetting (Mazurkiewicz) erasure folding exactly the commuting concurrent steps. Attaching
the virtual-token energy of the companion note as a $1$-cochain on the reduction graph, we find
that conservation of energy inverts the naive expectation: it is not bought by coarser erasure but
is a \emph{lower bound on how much history must be kept}, set by the arbitrage (holonomy) class of
the energy form. Frictionless (exact) energy is conserved at every grade, down to total erasure;
dissipative energy carries holonomy, and conserving it forces retention of exactly the loops that
holonomy obstructs---the holonomy-trivialising cover. Underneath sits a ledger law needing no
cohomology---initial fuel equals fuel remaining plus fuel recorded spent, a potential-plus-kinetic
invariant---whose one-way breakage under erasure is Landauer's principle: the history one may erase
for free is exactly the redundant, state-recoverable (exact) part, so the Hodge threshold of the
companion note and the Landauer threshold index the same class. Reversibility finally supplies the
time-reversal symmetry whose Noether charge is the conserved energy, offering the generative half
that the companion note's energy--time conjecture left open. Two structural refinements---the
sublattice of energy-conserving grades and the dagger structure of the free history---are noted and
deferred.
\end{abstract}

\section{Introduction}

The cost endofunctor $\cost$ on $\ciGSLT$ \cite{CE} meters a rewrite theory by gating each step on
the consumption of a token; the companion note \cite{VT} adjoins engines between token colours and
extracts a virtual token---an energy---as the potential whose existence is no-arbitrage. Both
constructions run \emph{forward}: they spend. This note adds the complementary gesture. A rewrite
theory is irreversible because reduction forgets: from $Q$ one cannot in general recover the $P$
that reduced to it, since many $P$ may. If instead the state carries a \emph{history}---a record of
the steps taken---and each rewrite \emph{appends} to it, then the predecessor is recorded and
reduction becomes invertible. Reversibility is not a new dynamics; it is the old dynamics run over
an enriched state that refuses to forget.

We make ``refuse to forget'' functorial. Section~\ref{sec:monad} presents the history endofunctor
$\hist$ as the writer monad over the free monoid of rewrite events, dual in character to $\cost$
and equipped with the same free--forgetful resolution, with the augmentation that drops history
being the reintroduction of irreversibility. Section~\ref{sec:tree} identifies the free history's
state graph as the universal cover of the reduction graph---a tree---and reversibility as its
unique-parent property. Section~\ref{sec:erasure} organises the erasures as a lattice of quotients
between the reversible cover and the irreversible base, distinguishing the erasures that merely
compress the log from those that fold the tree and re-admit loops. Section~\ref{sec:energy} carries
the virtual-token energy onto the reduction graph as a $1$-cochain, and Section~\ref{sec:price}
draws the note's main consequence: energy conservation is a lower bound on retained history, set by
the holonomy of that cochain. Section~\ref{sec:ledger} exhibits the underlying ledger law and reads
its one-way breakage as Landauer's principle, identifying the free-to-erase history with the exact
part of the energy form. Section~\ref{sec:noether} observes that reversibility supplies the
time-reversal symmetry whose Noether charge is the conserved energy, addressing the generative half
of the companion note's conjecture. Section~\ref{sec:further} records two deferred refinements.

Throughout, the register is that of a working note: several statements are given at the level of
rigour the setting affords, with the load-bearing refinements flagged as such.

\section{The history monad}
\label{sec:monad}

Fix a continued interactive GSLT $G=(\Sigma,E,R)$ with interacting sort $P$ and rewrite rules $R$.

\begin{definition}[Rewrite events]\label{def:event}
A \emph{rewrite event} for $G$ records everything needed to invert one step: the rule
$\ell\to r\in R$ fired, the occurrence (context) at which it fired, and, in a non-left-linear
theory, the material discarded or duplicated by the step together with the choice made among
concurrently enabled redexes. Write $\Ev$ for the set of events and
$M=(\Ev^{*},\cdot,\emptyhist)$ for the free monoid on $\Ev$: finite sequences of events under
concatenation, with the empty history $\emptyhist$.
\end{definition}

\begin{definition}[The history endofunctor]\label{def:hist}
$\hist(G)$ is the continued interactive GSLT whose interacting configurations are pairs $(P,h)$ with
$h\in M$, and whose rules are the \emph{logged} rules: for each $\ell\to r$ in $R$,
\[
  (\ell,\,h)\ \longrightarrow\ (r,\,h\cdot e_{\ell\to r}),
\]
appending the event $e_{\ell\to r}$ that names this step. Structural equations, the section, and
the token apparatus (if $G$ is already cost-accounted) are carried unchanged on the first
component; the history component is inert under $\equiv$.
\end{definition}

\begin{proposition}[$\hist$ is a monad, dual to $\cost$]\label{prop:monad}
$\hist$ carries the structure of a monad $(\hist,\eta,\mu)$ on $\ciGSLT$. The unit
$\eta_G:G\to\hist(G)$, $P\mapsto(P,\emptyhist)$, starts every term with the empty history; the
multiplication $\mu_G:\hist^2(G)\to\hist(G)$ concatenates a history-of-histories into a single
history. The augmentation $\pi_G:\hist(G)\to G$, $(P,h)\mapsto P$, forgets the history; it is a
$\ciGSLT$-morphism, and $\pi_G\circ\eta_G=\id_G$. As with $\cost$ \cite[\S9]{CE}, $\hist$ resolves
through a free--forgetful adjunction $\mathrm{Log}\dashv\mathrm{Forget}$ whose induced monad is
$\hist$: $\mathrm{Log}$ installs the history apparatus and $\mathrm{Forget}$ strips it, the
stripping being behaviour-altering---it reintroduces irreversibility.
\end{proposition}

\begin{remark}[Two ledgers of one clock]\label{rmk:duality}
$\hist$ and $\cost$ are the two temporal halves of a single construction. The token stack $::$ of
$\cost$ is a \emph{forward} modulus: it is produced up front and \emph{drained} by reduction, its
consumed prefix being the run length \cite[Prop.~4.1]{CE}. The history $h$ of $\hist$ is the
\emph{same quantity accrued from the other end}: it is produced by reduction and never drained, its
length being the run length again. One is the ledger of fuel remaining, the other the ledger of
fuel spent; together they book-keep one clock from both ends. This duality is sharpened into a
conservation law in Section~\ref{sec:ledger}, where the two lengths are shown to sum to a constant.
\end{remark}

\section{The free history is the reversible tree}
\label{sec:tree}

Let $\Grph(P)$ be the \emph{reduction graph} rooted at a term $P$: vertices are the configurations
reachable from $P$ under $R$, edges are single steps. Irreversibility of $G$ is visible here as the
failure of a unique-predecessor property---a vertex may have several in-edges (two terms reducing to
a common one) or lie on a self-loop (e.g.\ $\Omega\to\Omega$).

\begin{proposition}[Free history unfolds the reduction graph]\label{prop:cover}
Let $\hist_{\mathrm{free}}$ be the history functor of Definition~\ref{def:hist} with the full free
monoid $M$ (no erasure). The sub-LTS of $\hist_{\mathrm{free}}(G)$ reachable from
$(P,\emptyhist)$ is isomorphic to the tree unfolding of $\Grph(P)$---equivalently, to its universal
cover $\cover(P)$---with the augmentation $\pi_G$ as the covering map. Concretely, its vertices are
the finite reduction \emph{paths} out of $P$, an edge extends a path by one step, and the history
of a node is its address: the path from the root.
\end{proposition}

\begin{proof}[Proof sketch]
A reachable configuration $(Q,h)$ is determined by $h$, since $h$ is the sequence of steps taken
and $P$ is fixed; and every $h$ that is a valid step-sequence from $P$ occurs, ending at the term
$Q=\pi_G(Q,h)$ it induces. Thus reachable configurations are in bijection with finite reduction
paths from $P$, edges extend by one event, and no two distinct paths are identified---this is the
tree unfolding, whose defining universal property among covers of $\Grph(P)$ makes it the universal
cover; $\pi_G$ sends a path to its endpoint, a covering map.
\end{proof}

\begin{corollary}[History is reversibility]\label{cor:rev}
In $\hist_{\mathrm{free}}(G)$ every non-root node has a unique parent---delete the last event---so
the backward relation is a total function on the reachable tree: $\hist_{\mathrm{free}}(G)$ is
backward-deterministic. ``Keep a history'' and ``the computation graph is a tree'' are one
statement; the tree is the reversible cover, and the base $\Grph(P)$ is recovered by folding it
along $\pi_G$.
\end{corollary}

Corollary~\ref{cor:rev} is the abstract form of the constructions that make specific process
calculi reversible by remembering their past---CCSK \cite{CCSK} and RCCS \cite{RCCS}; here the
remembering is the free monad $\hist_{\mathrm{free}}$ and backward-determinism is the tree's
unique-parent property, before any calculus-specific bookkeeping.

\section{Erasures: the grading lattice}
\label{sec:erasure}

Between the reversible cover and the irreversible base lie the partial histories. An \emph{erasure}
is a congruence $\sim$ on the history monoid $M$ that is compatible with appending and with the
dynamics, yielding a quotient monad $\hist_{\sim}$ and a factorisation of the augmentation
\[
  \hist_{\mathrm{free}}(G)\ \twoheadrightarrow\ \hist_{\sim}(G)\ \twoheadrightarrow\ \id(G)=G,
\]
the right-hand terminal quotient collapsing all history. Under refinement these erasures form a
lattice, with $\hist_{\mathrm{free}}$ at the top (remember everything, fully reversible) and the
base at the bottom (remember nothing, fully irreversible).

Two qualitatively different bands of the lattice matter.

\begin{definition}[Log-compressing vs.\ folding erasures]\label{def:bands}
An erasure is \emph{log-compressing} if the induced map on reachable configurations remains
injective---distinct nodes of $\cover(P)$ stay distinct---so only the \emph{addresses} shorten while
the tree, and hence reversibility, is preserved. An erasure is \emph{folding} if it identifies
distinct nodes lying over a common base vertex, closing a loop; thereafter a path may revisit a
state and the quotient is a proper intermediate cover of $\Grph(P)$.
\end{definition}

\begin{proposition}[Folding erasures are intermediate covers]\label{prop:galois}
The folding erasures of $\hist_{\mathrm{free}}(G)$ over $\Grph(P)$ correspond, by the Galois theory
of covering spaces, to the subgroups of the fundamental group $\fund(\Grph(P))$---a free group of
rank $\Betti(\Grph(P))$. A folding erasure closes exactly the loops in its chosen subgroup;
choosing all loops recovers the base, choosing none recovers the tree.
\end{proposition}

\begin{remark}[The order-forgetting erasure is true concurrency]\label{rmk:mazurkiewicz}
One folding grade is distinguished. Abelianising the history---retaining the \emph{multiset} of
events rather than their order---identifies two paths iff they differ by transposing independent,
concurrently-enabled steps. This is the maximal abelian cover, and in a concurrent theory it is
exactly the Mazurkiewicz-trace \cite{Mazurkiewicz} (true-concurrency) unfolding: the grade at which
``same computation up to scheduling'' becomes ``same state.'' The purely sequential interleavings
are folded away; only genuinely causal branching remains in the tree.
\end{remark}

\begin{remark}[Loop rank here, arbitrage rank there]
The rank $\Betti(\Grph(P))$ governing the folding lattice is the first Betti number of the
\emph{reduction} graph, the exact analogue of the first Betti number of the \emph{engine} graph
that governed the arbitrage lattice of \cite[Prop.~4.3]{VT}. In both notes the room for a
degeneracy---loops to close here, arbitrage cycles to price there---is a first homology of a graph,
and the two constructions meet in the next two sections, where the energy of \cite{VT} is evaluated
along the loops of $\Grph(P)$.
\end{remark}

\section{Energy on the reduction graph}
\label{sec:energy}

The companion note \cite{VT} produces, from an arbitrage-free system of engines, a virtual-token
valuation $\nu$ assigning each token colour a real value, and---under friction---a dissipation
$\theta\ge 0$ charged by each lossy conversion \cite[\S7]{VT}. Transport this onto the reduction
graph as the energy increment effected by each step.

\begin{definition}[Energy $1$-cochain]\label{def:omega}
The \emph{energy form} $\omega\in C^1(\Grph(P);\Real)$ assigns to each step $Q\to Q'$ the
virtual-token energy it effects, inclusive of any dissipation: $\omega(Q\to Q')$ is the change in
$\nu$-value of the configuration's token content across the step, plus the dissipation $\theta$ of
any lossy engine the step invokes.
\end{definition}

By the graph--Hodge decomposition \cite{Hodge} used in \cite[\S4]{VT}, $\omega$ splits orthogonally,
\[
  \omega \;=\; \underbrace{\delta\phi}_{\text{exact}} \;\oplus\; \underbrace{\omega_{\mathrm{cyc}}}_{\text{cyclic}},
\]
with $\delta\phi$ the gradient of a state-energy potential $\phi$ on configurations and
$\omega_{\mathrm{cyc}}$ a divergence-free circulation. The holonomy of a loop $\gamma$ in
$\Grph(P)$ is
\[
  \hol(\gamma)\;=\;\oint_{\gamma}\omega\;=\;\oint_{\gamma}\omega_{\mathrm{cyc}},
\]
the exact part contributing nothing around closed paths.

\begin{proposition}[Exactness is frictionless energy]\label{prop:exact}
$\omega$ is exact---$\omega_{\mathrm{cyc}}=0$, all holonomies vanish---iff energy is a state
function: the value $\phi(Q)$ of a configuration is path-independent, so no loop of computation
changes the energy. This is the frictionless limit $\theta\equiv 0$. A nonzero holonomy
$\hol(\gamma)\neq 0$ is exactly a loop around which the dissipation $\oint_\gamma\theta$ fails to
cancel: energy tracked inclusive of loss is then path-dependent.
\end{proposition}

\section{The memory price of conservation}
\label{sec:price}

We can now state the note's main consequence. Fold the tree by some erasure; the folded loops are
those the erasure closes (Proposition~\ref{prop:galois}). Energy descends to a well-defined state
function on the folded quotient iff $\omega$ restricted to those loops is exact---iff every closed
loop has vanishing holonomy.

\begin{theorem}[Conservation bounds erasure from below]\label{thm:price}
Let $\sim$ be a folding erasure closing the loop-subgroup $\Lambda\le\fund(\Grph(P))$. Energy is
conserved on $\hist_{\sim}(G)$---i.e.\ $\phi$ descends to a single-valued state function on the
folded quotient---iff $\hol(\gamma)=0$ for every $\gamma\in\Lambda$. Consequently:
\begin{enumerate}[label=\textup{(\roman*)}]
\item If $\omega$ is exact (arbitrage-free, frictionless), every $\Lambda$ satisfies the condition:
energy is conserved at \emph{every} grade, down to and including total erasure. A conservative
system tolerates arbitrary forgetting.
\item If $\omega$ carries a nonzero holonomy class, the grades on which energy is conserved are
exactly those whose closed loops avoid it. Conservation then \emph{forbids} folding the
holonomy-carrying loops: they must remain unfolded as tree-branches, i.e.\ history recording them
must be retained.
\end{enumerate}
\end{theorem}

\begin{proof}[Proof sketch]
A potential descends along a folding iff it is constant on each identified fibre, i.e.\ iff the
increment around each closed loop is zero; that increment is $\hol(\gamma)$. (i) and (ii) are the
two cases $\omega_{\mathrm{cyc}}=0$ and $\omega_{\mathrm{cyc}}\neq0$.
\end{proof}

The naive expectation---that coarser erasure buys conservation---is thus inverted. Conservation is
a \emph{lower bound on how much history must be kept}, and the bound is set by the holonomy class of
$\omega$: zero holonomy puts the bound on the floor (erase everything and energy survives), while
nonzero holonomy forces retention up to the cover that trivialises it. That cover is the
Riemann-surface move---unwinding a multivalued potential by climbing onto the sheet where it becomes
single-valued, remembering the winding rather than collapsing it.

\begin{corollary}[The maximal energy-conserving grade]\label{cor:maximal}
Among the erasures on which energy is conserved there is a coarsest---fold the tree everywhere the
potential is single-valued and leave it unfolded exactly where the holonomy obstructs, i.e.\ the
cover corresponding to the smallest normal subgroup of $\fund(\Grph(P))$ containing every
zero-holonomy loop. This is the natural physical state space of the computation: the reduction
graph quotiented as far as energy conservation permits and no further.
\end{corollary}

We record that the existence of Corollary~\ref{cor:maximal}'s coarsest grade in full generality---as
opposed to under finiteness of $\Grph(P)$---depends on the erasure lattice possessing the relevant
joins, and defer this to Section~\ref{sec:further}.

\section{The ledger law and Landauer's principle}
\label{sec:ledger}

Beneath the cohomology sits an elementary invariant, and it is the cleanest thing to state.

\begin{proposition}[The ledger law: potential plus kinetic]\label{prop:ledger}
Run a cost-accounted term under $\hist$ from an initial stack of depth $\sigma_0$ (total fuel). Let
$\sigma(t)$ be the stack depth remaining and $\kappa(t)$ the number of events recorded in the
history after $t$ steps. Each forced step pops one cell and appends one event, so
\[
  \sigma(t)+\kappa(t)\;=\;\sigma_0\qquad\text{for all }t.
\]
Fuel remaining plus fuel recorded spent is conserved by construction.
\end{proposition}

Read $\sigma$ as potential energy---fuel not yet expended---and $\kappa$ as kinetic---work done and
logged. Their sum is the conserved total. Reversibility is that the transfer is invertible: undo a
step (pop an event, push a cell) and one unit returns from spent to remaining. This is the
frictionless, exact case of Section~\ref{sec:energy} read on the nose.

Erasure is what makes the transfer one-way. Forget a recorded event and $\kappa$ falls, but $\sigma$
cannot rise---one cannot un-spend fuel---so the invariant breaks \emph{irreversibly}.

\begin{remark}[Landauer's principle]\label{rmk:landauer}
The irreversible step is erasure, and erasure has a cost. Discarding a history event that is
\emph{not} recoverable from the current state destroys information and, by Landauer's principle
\cite{Landauer,Bennett}, dissipates energy as heat; discarding an event that \emph{is} recoverable
from the state is logically reversible and free. The recoverable events are precisely those
determined by the current configuration---the \emph{exact} part $\delta\phi$ of the energy form,
the forced reversible returns---while the irrecoverable events are the holonomy-carrying part
$\omega_{\mathrm{cyc}}$, the genuinely lost information. Hence the history one may erase for free is
exactly the history whose erasure preserves energy: \emph{erase-freely} and
\emph{conserve-energy-freely} are the same permission, bounded by the same class. The Hodge
threshold of \cite{VT} and the Landauer threshold index one and the same cohomology class of
$\omega$.
\end{remark}

Remark~\ref{rmk:landauer} is the mechanism under the exergy reading of \cite[\S7]{VT}: dissipation
is the Landauer heat of erasing the irrecoverable, holonomy-carrying history, and the free energy
that decreases monotonically is what remains recoverable.

\section{Reversal symmetry and the generative half of energy}
\label{sec:noether}

The companion note closed with an energy--time conjecture \cite[Conj.~8.1]{VT}: the virtual token is
manifestly the conserved \emph{charge}, but the token stack alone---one-way time, only draining---does
not exhibit it as the \emph{generator} of the dynamics that physical energy (the Hamiltonian) is.
The history functor supplies the missing direction.

\begin{conjecture}[Reversal symmetry carries the generator]\label{conj:noether}
In $\hist_{\mathrm{free}}(G)$ reduction is time-reversal symmetric: each forward step has a unique
inverse (Corollary~\ref{cor:rev}), so the augmented dynamics admits an involution exchanging
forward reduction with backward replay. This involution is the time-reversal symmetry whose Noether
\cite{Noether} charge is the conserved energy of Theorem~\ref{thm:price}(i); the two-sided time of
Remark~\ref{rmk:duality}---$::$ forward, $h$ backward, balanced by the ledger law---is the parameter
it translates, and the conserved energy is its generator. The erasure grades measure how much of the
symmetry, hence how much of the conservation, is spent: total erasure breaks the involution
entirely and returns the one-way, non-generative charge of the stack alone.
\end{conjecture}

Conjecture~\ref{conj:noether} is offered as the reversible completion of \cite[Conj.~8.1]{VT}: the
stack gives energy as a debit, the free history restores the reverse direction and hence the
symmetry, and Noether's correspondence---where it can be made to hold---turns the symmetry's charge
into the generator. Whether the involution is literally a dagger on the augmented LTS is the
subject of the next section.

\section{Two deferred refinements}
\label{sec:further}

Two structural questions are load-bearing and left for a subsequent note.

\emph{The sublattice of energy-conserving grades.} Whether the erasures compatible with the rewrite
structure form a lattice possessing the joins needed for the coarsest energy-conserving grade of
Corollary~\ref{cor:maximal} to exist in general, or only under finiteness of $\Grph(P)$, is open.
The covering-space correspondence of Proposition~\ref{prop:galois} suggests the answer is the
lattice of subgroups of $\fund(\Grph(P))$ cut down to those closing only zero-holonomy loops; making
this a genuine sublattice with a top element below the tree is the first refinement.

\emph{The dagger structure of the free history.} Backward-determinism (Corollary~\ref{cor:rev})
gives an involution on the reachable tree, but whether it is a genuine dagger on the augmented LTS
of a \emph{non-linear} GSLT---where the discarded-and-duplicated-material bookkeeping of
Definition~\ref{def:event} must be inverted exactly---is where a clean involution could fail. The
reversible process calculi \cite{CCSK,RCCS} obtain backward-determinism in the linear setting; the
non-linear case, and hence the precise sense of Conjecture~\ref{conj:noether}, is the second
refinement.

\section*{Conclusion}
\addcontentsline{toc}{section}{Conclusion}

Keeping a history and appending at every rewrite is an endofunctor $\hist$ on $\ciGSLT$, the writer
monad over rewrite events and the time-reversal dual of the cost endofunctor: $\cost$ drains a
forward modulus, $\hist$ accrues a backward one, and the ledger law binds the two into a conserved
sum. The free history is the universal cover of the reduction graph---a tree---and reversibility is
its unique-parent property; erasures grade the passage from this reversible cover down to the
irreversible base through the intermediate covers, the order-forgetting grade landing on true
concurrency. Evaluating the virtual-token energy of the companion note along the loops of the
reduction graph inverts the naive picture of conservation: energy is conserved not by forgetting
more but by remembering enough, the memory bound set by the holonomy class of the energy form, and
the history free to erase is exactly the state-recoverable, exact part---so that the Hodge threshold
of the companion note and Landauer's threshold coincide. Reversibility supplies the time-reversal
symmetry whose Noether charge is the conserved energy, completing from the reverse direction the
energy--time picture the companion note began. The construction is thus the temporal complement of
cost accounting: where metering asks what a computation may spend, history asks what it must
remember, and energy conservation is exactly the exchange rate between the two.

\bigskip
\noindent\emph{Nihil Sine Vestigio.}

\begin{thebibliography}{99}

\bibitem{VT} L.~G.~Meredith, \emph{Energy as a Virtual Token: No-Arbitrage, the Graph--Hodge
Potential, and the Conserved Scalar of Cost-Accounted GSLTs}, F1R3FLY.io working note, 2026.

\bibitem{CE} L.~G.~Meredith, \emph{Continued Interactive GSLTs and the Cost Endofunctor: A
construction one level up from cost-accounted Rholang}, F1R3FLY.io, 2026.

\bibitem{CARho} L.~G.~Meredith, \emph{Cost-Accounted Rho Calculus: A Spectral Decomposition of
Phlogiston}, F1R3FLY.io, 2026.

\bibitem{MC} L.~G.~Meredith, \emph{Mortal Computation: Mortality, the Genome, and Life-History in
Cost-Accounted GSLTs}, F1R3FLY.io working note, 2026.

\bibitem{CCSK} I.~Phillips and I.~Ulidowski, \emph{Reversing algebraic process calculi}, Journal of
Logic and Algebraic Programming, vol.~73, no.~1--2, pp.~70--96, 2007.

\bibitem{RCCS} V.~Danos and J.~Krivine, \emph{Reversible communicating systems}, in CONCUR 2004,
Lecture Notes in Computer Science, vol.~3170, pp.~292--307, Springer, 2004.

\bibitem{Landauer} R.~Landauer, \emph{Irreversibility and heat generation in the computing process},
IBM Journal of Research and Development, vol.~5, no.~3, pp.~183--191, 1961.

\bibitem{Bennett} C.~H.~Bennett, \emph{Logical reversibility of computation}, IBM Journal of
Research and Development, vol.~17, no.~6, pp.~525--532, 1973.

\bibitem{Mazurkiewicz} A.~Mazurkiewicz, \emph{Trace theory}, in Petri Nets: Applications and
Relationships to Other Models of Concurrency, Lecture Notes in Computer Science, vol.~255,
pp.~278--324, Springer, 1987.

\bibitem{Moggi} E.~Moggi, \emph{Notions of computation and monads}, Information and Computation,
vol.~93, no.~1, pp.~55--92, 1991.

\bibitem{Hodge} X.~Jiang, L.-H.~Lim, Y.~Yao, and Y.~Ye, \emph{Statistical ranking and combinatorial
Hodge theory}, Mathematical Programming, vol.~127, no.~1, pp.~203--244, 2011.

\bibitem{Noether} E.~Noether, \emph{Invariante Variationsprobleme}, Nachrichten von der
Gesellschaft der Wissenschaften zu G\"ottingen, Mathematisch-Physikalische Klasse, pp.~235--257,
1918.

\bibitem{MacLane} S.~Mac~Lane, \emph{Categories for the Working Mathematician}, 2nd ed., Graduate
Texts in Mathematics, vol.~5, Springer, 1998.

\end{thebibliography}

\end{document}
